\section{Standard tensor layers with an entire external partner}
\label{sec:standard-entire-partner}

\paragraph{Status and provenance.}
HD-AST13--15, reserved question Q2dx, is a separate derivation by
\texttt{/root/imprimitive} following the parent's proposed extension
to the entire external partner.  The accepted HD-AST07 pair inequality
and HD-TW04 weighted-cycle summation are shared prior inputs.  The new
point is to use two different local inequalities: the stronger pair
bound at factor-fixing entries and a weaker single-factor bound around
arbitrary moving cycles.  The complete author proof, code and output
readings pass. The complete root and independent whole-proof, source and code readings also pass without mathematical repair.  It leaves every
previously accepted AST and TW section unchanged.

Let $n\ge15$, let $W=\operatorname{Std}_n$ be the ordinary deleted
permutation module of $S_n$, and set $a=n-1$.  Consider an actual
finite coefficient image on
\[
 V=W^{\otimes t}\otimes U,\qquad t\ge2,\quad b=\dim U\ge1,
 \qquad e=a^t b.
 \tag{ASTP.1}
\]
The first tensor factor is stable as a whole.  Its actual projective
image contains the represented ordinary core $A_n^t$, normalizes that
core, and permutes the $t$ component algebras transitively.  The partner
has its actual finite projective factor action and satisfies
\[
 \frac{\nu_U(S)}b\ge\frac1n
 \quad\hbox{for every nonscalar operator in that partner image},
 \qquad
 \nu(T)=\dim T-\max_z\dim\ker(T-zI).
 \tag{ASTP.2}
\]
For $b=1$ the condition is vacuous.  Every assertion permits arbitrary
scalar phases and all original coefficient fields.

\subsection*{HD-AST13. The two different local standard estimates}

\begin{lemma}[HD-AST13]
For every $x\in S_n$, with $k_x=\operatorname{ind}(x)$, one has
\[
 \nu(W(x))\ge\frac{n-2}{n}k_x
             \ge A k_x,
 \qquad A=\frac{n-3}{a}.
 \tag{ASTP.3}
\]
Put $B=n/(2a)$.  At every genuine coefficient inertia which fixes
all the $t$ standard factors, let $I$ be the sum of its within-factor
permutation indices.  Then, including its external partner and actual
scalar phase,
\[
                 \frac{\operatorname{rank}(1-T)}e
                   \ge\frac{A I}{ta}+\frac Ba.
 \tag{ASTP.4}
\]
The stronger expression $A k_x+B$ is not claimed for a single
standard factor.  In particular the proof does not substitute it
into the hypothesis (TWP.3) of HD-TW04.
\end{lemma}

\begin{proof}
Write $c=n-k_x$ for the number of cycles, including fixed points.
On the full letter module every target eigenspace has dimension at
most $c$, since each cycle contributes at most one.  On the deleted
module the eigenvalue-one multiplicity is $c-1$.
If $c>n/2$, a nonidentity eigenvalue cannot receive a contribution
from all $c$ cycles: that would force every cycle to have length at
least two, contradicting $2c>n$.  Its multiplicity is therefore
at most $c-1$ as well.  In this case $\nu(W(x))\ge k_x$.
If $c\le n/2$, the general multiplicity bound gives
$\nu(W(x))\ge a-c=k_x-1$, and $k_x\ge n/2$ implies
\[
 k_x-1\ge\frac{n-2}{n}k_x,
 \qquad
 \frac{n-2}{n}-\frac{n-3}{n-1}=\frac2{n(n-1)}>0.
 \tag{ASTP.5}
\]
This proves (ASTP.3), including the identity case separately with
both sides zero and all off-spectrum targets with full moving rank.

For (ASTP.4), suppose $I>0$.  The average of the index sums over
the unordered pairs of tensor factors is $2I/t$.  Choose a pair
with at least this sum.  It is not the identity pair.  The already
proved all-target HD-AST07 inequality gives projective support on
this pair at least
\[
 \frac{(n-3)(k_x+k_y)+n}{2}
       \ge\frac{n-3}{t}I+\frac n2.
 \tag{ASTP.6}
\]
Tensor spectral convolution with every other standard factor and
with the arbitrary external partner preserves the normalized
support.  Dividing (ASTP.6) by $a^2$ is exactly (ASTP.4).
This is the pair estimate; it does not rely on (ASTP.3) to supply
its positive constant term.

If $I=0$, all constituent permutations are the identity, so the
whole first factor is scalar.  A nonscalar external operator has
normalized moving rank at least $1/n$ by (ASTP.2), and
\[
              \frac1n\ge\frac{n}{2a^2}=\frac Ba
              \qquad(n\ge15).
 \tag{ASTP.7}
\]
Indeed $2(n-1)^2-n^2=n^2-4n+2>0$ in this range.
If the partner is scalar as well, genuineness of the full inertia
means that the total scalar is not one; its normalized moving
rank is one, again at least $B/a$.  This proves every case of
(ASTP.4).  Thus the old scalar-only identity-projection step has
been replaced by the actual whole-partner floor.
\end{proof}

\subsection*{HD-AST14. Exclusion at every tensor count at least two}

\begin{theorem}[HD-AST14]
There is no connected finite-group sphere cover in the setting
(ASTP.1)--(ASTP.2) with a generating product-one tuple of $r\ge12$
genuine coefficient inertias and
\[
                    \sum_{i=1}^r\operatorname{rank}(1-T_i)=2e+6.
 \tag{ASTP.8}
\]
Consequently such a standard tensor layer cannot occur in an actual
rank-nine sixfold supplier after the usual intermediate-cover
reduction.  This statement holds over every original imaginary
quadratic field, including the Eisenstein field.  It uses no scalar
conductor marks.
\end{theorem}

\begin{proof}
We first verify the actual quotient interface in the presence of $U$.
The represented perfect core $E=A_n^t$ acts irreducibly on
$W^{\otimes t}$.  Projective normalization already gives linear
normalization: if $C E C^{-1}\subset\mathbf C^*E$, taking derived
groups gives $C E C^{-1}\subset E$, and applying this to $C^{-1}$
gives equality.  The accepted compatible-normalizer calculation
HD-AST04, using $\operatorname{Aut}(A_n)=S_n$ for $n\ge7$, says
that this normalizer is the product of the scalars with the rational
tensor-wreath image of $S_n\wr S_t$.

The rational tensor-wreath image intersects the scalars trivially.
Indeed a scalar centralizes each individual $A_n$ component, hence
has trivial factor permutation; each remaining $S_n$ element then
centralizes its $A_n$ and is the identity.  The same reasoning
proves faithfulness of that rational wreath representation.  Thus
the actual projective first-factor action has a unique homomorphic
lift to this rational wreath group.  In particular its actions on
$t$ factors and on $nt$ letters are actual quotient actions of the
given coefficient cover, not branchwise choices.  The external
partner and its scalar choices do not change these homomorphisms.
The $nt$-letter action is transitive because the image contains the
ordinary core, transitive within each block, and is transitive on
the blocks by hypothesis.

A factor-moving entry has a weighted tensor cycle of length
$\ell\ge2$.  In an eigenbasis for its finite-projective-order
cycle product, $a$ basis words are constant and every other rotation
orbit has length at least two.  Each weighted orbit contributes
at most one vector to a specified eigenspace.  Its normalized
support is consequently at least $(a-1)/(2a)$; convolution with
the other cycles and $U$ preserves that bound.  Since
$(a-1)/(2a)\ge B/a$ for $n\ge15$, HD-AST13 implies that every
genuine entry, moving or fixed, costs at least $B/a$.

If at least four entries move factors, select four such entries
and eight further genuine entries.  Their moving ranks imply
\[
 \begin{split}
 \sum_i\operatorname{rank}(1-T_i)
 &\ge4\frac{a-1}{2a}e+8\frac Ba e\\
 &=2e+\frac{e}{a^2}(2n+2)>2e+6,
 \end{split}
 \tag{ASTP.9}
\]
because $e/a^2=a^{t-2}b\ge1$.  Thus at most three entries move
factors, and at least nine fix them.  No assertion that precisely
two move is needed.

For a general entry let $\pi_i$ be its factor permutation, with
cycles of lengths $\ell_j$ and compatible standard cycle products
$x_j\in S_n$.  Let $I_i$ denote its actual $nt$-letter index.
The exact cycle-index identity and the accepted weighted-cycle
bound are
\[
 \begin{gathered}
 I_i=n\operatorname{ind}(\pi_i)+\sum_j k_{x_j},\\
 s_j\ge1-\frac1{\ell_j}
      +\frac{\nu(W(x_j))}{\ell_j a}
      -\frac{R_{\ell_j}(a)}{\ell_j a^{\ell_j}},\qquad
 R_1=0,\quad \frac{R_\ell(a)}{a^\ell}\le\frac1a\ (\ell\ge2).
 \end{gathered}
 \tag{ASTP.10}
\]
Here $R_\ell=\sum_{j=1}^{\ell-1}a^{\gcd(\ell,j)}$ is the exact
short-word error of D-ATT08/HD-AST11; its uniform bound for
$a\ge4$ and all $\ell\ge2$ is the accepted HD-TW01 input.
Convolution bounds the full support below by the weighted average
$\sum_j(\ell_j/t)s_j$.  There are at most
$\operatorname{ind}(\pi_i)$ nontrivial cycles.  Using only the
weaker single-factor inequality (ASTP.3), including $x_j=1$, gives
\[
 \frac{\operatorname{rank}(1-T_i)}e
 \ge\frac{A I_i}{ta}
       +\frac{\lambda}{t}\operatorname{ind}(\pi_i),
 \quad
 \lambda=1-\frac{nA+1}{a}=\frac2{a^2}>0,
 \quad L=a-(n-1)A=2.
 \tag{ASTP.11}
\]
The extra $B/a$ from (ASTP.4) is available at every factor-fixing
entry.  The two connected actual quotient covers give
\[
          \sum_i I_i\ge2nt-2,
          \qquad\sum_i\operatorname{ind}(\pi_i)\ge2t-2.
 \tag{ASTP.12}
\]
Original marks whose projection is identity contribute zero to
these index sums but are still among the genuine coefficient
entries charged in (ASTP.4).

Because $A\ge0$, $\lambda>0$, and $r_{\rm fix}\ge9$, summation
of (ASTP.11)--(ASTP.12) and the fixed-entry constants gives
\[
 \begin{split}
 \sum_i\operatorname{rank}(1-T_i)-2e
 &\ge\frac ea\left(r_{\rm fix}B-2-\frac2t(L-1)\right)\\
 &\ge\frac ea(9B-3)
   =\frac{e}{a^2}\left(\frac{3n}{2}+3\right)>6.
 \end{split}
 \tag{ASTP.13}
\]
This contradicts (ASTP.8).  The stated rank-nine consequence uses
the previously proved actual intermediate-cover factorization:
for $e>6$, the multiplicity Euler identity excludes base genus
at least two, and genus one has only six units of moving rank,
hence cannot have period rank nine.  On the sphere a rank-nine
image needs at least twelve genuine marks, giving exactly (ASTP.8).
No dimension inequality is being used as a sufficiency statement.
\end{proof}

\subsection*{HD-AST15. A maximal standard layer in a complete primitive decomposition}

\begin{corollary}[HD-AST15]
Let the actual coefficient image be finite and primitive, and use
the complete recursive decomposition D-TD12.  Suppose one extracted
nonabelian layer consists of $t\ge2$ ordinary standard components
with simple quotient $A_n$, where $n\ge80$, and $n$ is maximal
among all alternating natural degrees in every layer of this chosen
complete decomposition.  Then it cannot occur in a rank-nine
sixfold curve supplier, over any original imaginary quadratic field.
\end{corollary}

\begin{proof}
At this extraction layer the components are permuted transitively,
the whole first tensor factor is stable, and its actual projective
image contains their ordinary core.  D-TD12 gives a genuine
projective quotient on the layer and pulls it back through all
preceding multiplicity quotients.  Its external partner consists
of all the remaining extracted pieces.  Their complete family of
factor algebras is still normalized or permuted by the original
image; none is discarded.  The local cores in that partner have
alternating natural degrees at most $n$ by maximality.  Therefore
HD-FT21, using the completeness established by D-TD12 for this
exhibited factorization, gives
\[
 \frac{\nu_U(S)}{\dim U}
       \ge\frac1{\max\{80,m_{*,U}-1\}}\ge\frac1n.
 \tag{ASTP.14}
\]
For an empty partner the floor is vacuous.  All hypotheses of
HD-AST14 follow, so that theorem supplies the contradiction.
This argument uses the entire transitive layer, never an invented
stable individual $A_n$ quotient of a permuted layer.
\end{proof}

\paragraph{Exact remaining scope.}
The result treats every $t\ge2$ and genuinely extends the earlier
pure-tensor standard argument to arbitrary entire partners satisfying
(ASTP.2).  It also avoids the old non-Eisenstein conductor condition.
The case $t=1$ is explicitly outside these theorems.  If $n<80$, the
universal complete-partner floor alone need not imply $1/n$; the
main theorem still applies when that stronger actual floor is
separately established.  No bounded original induction index, field
descent, polarization compatibility, or existence of a supplier is
inferred from this necessary source obstruction.
